\documentclass[10pt]{extarticle}
\usepackage[a4paper,margin=1.5cm]{geometry}
\usepackage[T1]{fontenc}
\usepackage{amsmath,amssymb}
\usepackage{graphicx}
\usepackage[export]{adjustbox}
\usepackage{multicol}
\usepackage{enumitem}
\usepackage{booktabs,array,tabularx}
\setlength{\columnsep}{0.7cm}
\setlength{\parindent}{0pt}
\setlength{\parskip}{3pt}
\pagestyle{empty}
\makeatletter
% figure registry (figdefs.tex is generated by the extract step)
\newcommand{\qfigdef}[4]{\@namedef{qf@file@#1}{#2}\@namedef{qf@w@#1}{#3}\@namedef{qf@wide@#1}{#4}}
\newcommand{\qf@check}[1]{\@ifundefined{qf@file@#1}{\PackageError{score2}{Unknown figure id #1}{}}{}}
\newcommand{\qf@img}[4][]{\includegraphics[width=#2\dimexpr\@nameuse{qf@w@#4}pt\relax,max width=#3,max height=0.45\textheight,#1]{\@nameuse{qf@file@#4}}}
% Full-width questions: finalize wraps every question that contains a wide figure in \qfullbegin...\qfullend,
% so the WHOLE question (stem, figures, options) is set across the page in reading order, then the two columns
% resume. The question's head - its number, stem and first figure (\qfullsplit marks the end of the head) -
% is kept on one page, so a number is never left at a page foot with its figure overleaf; the rest may break
% across pages normally, so no big gaps are left. Without a \qfullsplit, the whole question is kept together
% if it is under 45% of a page. Never write these yourself.
\newif\ifqf@full
\newif\ifqf@boxed
\newsavebox{\qf@box}
\newcommand{\qf@open}{\global\qf@fulltrue\par\global\qf@boxedtrue
  \setbox\qf@box\vbox\bgroup\hsize\textwidth\linewidth\textwidth}
\newcommand{\qf@emit}[1]{\par\egroup\global\qf@boxedfalse
  \ifdim\dimexpr\ht\qf@box+\dp\qf@box\relax<#1\textheight\noindent\usebox{\qf@box}\par\else\unvbox\qf@box\fi}
\newcommand{\qfullsplit}{\ifqf@boxed\qf@emit{0.85}\fi}
\newcommand{\qf@close}{\ifqf@boxed\qf@emit{0.45}\fi\global\qf@fullfalse}
\newcommand{\qfullbegin}{\end{multicols}\qf@open}
\newcommand{\qfullend}{\qf@close\begin{multicols}{2}}
\newcommand{\qfullstart}{\qf@open}
\newcommand{\qfulldone}{\qf@close}
% \qfig[scale]{id}: block figure, centred, directly where it is written (size set by the tool).
% Fallback if a wide figure is not inside a full-width question: break out of the columns in place - no floats.
\newcommand{\qfig}[2][1]{\qf@check{#2}%
  \ifqf@full\qf@blk{#1}{#2}\else
  \ifnum\@nameuse{qf@wide@#2}=1 \expandafter\@firstoftwo\else\expandafter\@secondoftwo\fi
  {\end{multicols}\begin{center}\qf@img{#1}{\textwidth}{#2}\end{center}\begin{multicols}{2}}%
  {\qf@blk{#1}{#2}}\fi}
\newcommand{\qf@blk}[2]{\par\smallskip{\centering\qf@img{#1}{\linewidth}{#2}\par}\smallskip}
% \qfigi[scale]{id}: inline figure (inside option lists / tables / minipages); never breaks columns.
\newcommand{\qfigi}[2][1]{\qf@check{#2}\qf@img[valign=c]{#1}{\linewidth}{#2}}
\makeatother
% structure macros - lint reads these to check completeness and figure placement
\newcommand{\papersec}[1]{\par\bigskip{\centering\Large\bfseries PAPER-#1\par}\medskip}
\newcommand{\partsec}[2]{\par\medskip\noindent{\large\bfseries PART-#1: #2}\par\smallskip}
\newcommand{\sect}[1]{\par\smallskip\noindent\textbf{#1}\par}
\newcommand{\qn}[1]{\par\smallskip\noindent\textbf{#1.}~}
\newcommand{\paraq}[2]{\par\smallskip\noindent\textbf{Paragraph for Questions #1 and #2}\par}
\newenvironment{qoptions}{\begin{itemize}[leftmargin=1.8em,itemsep=1pt,topsep=2pt,label={}]}{\end{itemize}}
\newcommand{\opt}[2]{\item[(#1)] #2}
% figure registry (generated by score2_tools.py finalize)
\qfigdef{p1_phy_s1_q01_1}{Score Advanced - Test 15 (18 May 2023)-img/p1_phy_s1_q01_1.png}{107.2}{0}
\qfigdef{p1_phy_s1_q03_1}{Score Advanced - Test 15 (18 May 2023)-img/p1_phy_s1_q03_1.png}{82.7}{0}
\qfigdef{p1_phy_s1_q04_1}{Score Advanced - Test 15 (18 May 2023)-img/p1_phy_s1_q04_1.png}{99.7}{0}
\qfigdef{p1_phy_s1_q05_1}{Score Advanced - Test 15 (18 May 2023)-img/p1_phy_s1_q05_1.png}{134.0}{0}
\qfigdef{p1_phy_s2_q07_1}{Score Advanced - Test 15 (18 May 2023)-img/p1_phy_s2_q07_1.png}{121.0}{0}
\qfigdef{p1_phy_s2_q08_1}{Score Advanced - Test 15 (18 May 2023)-img/p1_phy_s2_q08_1.png}{52.7}{0}
\qfigdef{p1_phy_s2_q09_1}{Score Advanced - Test 15 (18 May 2023)-img/p1_phy_s2_q09_1.png}{72.0}{0}
\qfigdef{p1_phy_s2_q10_1}{Score Advanced - Test 15 (18 May 2023)-img/p1_phy_s2_q10_1.png}{298.8}{1}
\qfigdef{p1_phy_s2_q11_1}{Score Advanced - Test 15 (18 May 2023)-img/p1_phy_s2_q11_1.png}{76.5}{0}
\qfigdef{p1_phy_s2_q12_1}{Score Advanced - Test 15 (18 May 2023)-img/p1_phy_s2_q12_1.png}{131.0}{0}
\qfigdef{p1_phy_s3_q01_1}{Score Advanced - Test 15 (18 May 2023)-img/p1_phy_s3_q01_1.png}{101.4}{0}
\qfigdef{p1_phy_s3_q02_1}{Score Advanced - Test 15 (18 May 2023)-img/p1_phy_s3_q02_1.png}{125.6}{0}
\qfigdef{p1_phy_s3_q04_1}{Score Advanced - Test 15 (18 May 2023)-img/p1_phy_s3_q04_1.png}{112.9}{0}
\qfigdef{p1_phy_s3_q05_1}{Score Advanced - Test 15 (18 May 2023)-img/p1_phy_s3_q05_1.png}{104.9}{0}
\qfigdef{p1_phy_s3_q06_1}{Score Advanced - Test 15 (18 May 2023)-img/p1_phy_s3_q06_1.png}{117.3}{0}
\qfigdef{p1_chem_s1_q03_1}{Score Advanced - Test 15 (18 May 2023)-img/p1_chem_s1_q03_1.png}{80.1}{0}
\qfigdef{p1_chem_s1_q04_1}{Score Advanced - Test 15 (18 May 2023)-img/p1_chem_s1_q04_1.png}{308.0}{1}
\qfigdef{p1_chem_s1_q04_A}{Score Advanced - Test 15 (18 May 2023)-img/p1_chem_s1_q04_A.png}{73.4}{0}
\qfigdef{p1_chem_s1_q04_B}{Score Advanced - Test 15 (18 May 2023)-img/p1_chem_s1_q04_B.png}{64.1}{0}
\qfigdef{p1_chem_s1_q04_C}{Score Advanced - Test 15 (18 May 2023)-img/p1_chem_s1_q04_C.png}{52.3}{0}
\qfigdef{p1_chem_s1_q04_D}{Score Advanced - Test 15 (18 May 2023)-img/p1_chem_s1_q04_D.png}{53.7}{0}
\qfigdef{p1_chem_s2_q07_1}{Score Advanced - Test 15 (18 May 2023)-img/p1_chem_s2_q07_1.png}{73.2}{0}
\qfigdef{p1_chem_s2_q08_A}{Score Advanced - Test 15 (18 May 2023)-img/p1_chem_s2_q08_A.png}{91.1}{0}
\qfigdef{p1_chem_s2_q08_B}{Score Advanced - Test 15 (18 May 2023)-img/p1_chem_s2_q08_B.png}{90.3}{0}
\qfigdef{p1_chem_s2_q08_C}{Score Advanced - Test 15 (18 May 2023)-img/p1_chem_s2_q08_C.png}{88.9}{0}
\qfigdef{p1_chem_s2_q08_D}{Score Advanced - Test 15 (18 May 2023)-img/p1_chem_s2_q08_D.png}{92.5}{0}
\qfigdef{p1_chem_s2_q09_1}{Score Advanced - Test 15 (18 May 2023)-img/p1_chem_s2_q09_1.png}{280.7}{1}
\qfigdef{p1_chem_s2_q09_3}{Score Advanced - Test 15 (18 May 2023)-img/p1_chem_s2_q09_3.png}{58.0}{0}
\qfigdef{p1_chem_s2_q10_A}{Score Advanced - Test 15 (18 May 2023)-img/p1_chem_s2_q10_A.png}{29.9}{0}
\qfigdef{p1_chem_s2_q10_B}{Score Advanced - Test 15 (18 May 2023)-img/p1_chem_s2_q10_B.png}{31.3}{0}
\qfigdef{p1_chem_s2_q10_C}{Score Advanced - Test 15 (18 May 2023)-img/p1_chem_s2_q10_C.png}{51.4}{0}
\qfigdef{p1_chem_s2_q10_D}{Score Advanced - Test 15 (18 May 2023)-img/p1_chem_s2_q10_D.png}{48.9}{0}
\qfigdef{p1_chem_s3_q03_1}{Score Advanced - Test 15 (18 May 2023)-img/p1_chem_s3_q03_1.png}{163.8}{0}
\qfigdef{p1_chem_s3_q04_1}{Score Advanced - Test 15 (18 May 2023)-img/p1_chem_s3_q04_1.png}{289.0}{1}
\qfigdef{p1_chem_s3_q05_1}{Score Advanced - Test 15 (18 May 2023)-img/p1_chem_s3_q05_1.png}{202.3}{0}
\raggedcolumns
\makeatletter\@beginparpenalty=10000\makeatother
\newcolumntype{L}[1]{>{\raggedright\arraybackslash}p{#1}}
\begin{document}
\begin{center}{\Large\bfseries Score Advanced --- Test 15 (18 May 2023)}\end{center}
\begin{multicols}{2}
\raggedcolumns
\papersec{1}
\partsec{1}{PHYSICS}
\sect{SECTION-I (i)}
\qn{1} A taut string consists of two segments with linear mass densities $\mu$ kg/m to the left of A and $4\mu$ kg/m to the right of A. A sinusoidal pulse travels towards the right as shown.\par
\qfig{p1_phy_s1_q01_1}
\begin{qoptions}
\opt{A}{Ratio of reflected to transmitted amplitude = $\dfrac14$}
\opt{B}{Ratio of reflected to transmitted amplitude = 2}
\opt{C}{Ratio of transmitted to reflected power = 8}
\opt{D}{Ratio of transmitted to reflected power = 16}
\end{qoptions}
\qn{2} One mole of a monoatomic gas is heated in such a way that its molar specific heat is $2R$. If the volume of the gas is doubled, by what factor does the temperature change?
\begin{qoptions}
\opt{A}{$\dfrac{1}{\sqrt2}$}
\opt{B}{$\sqrt2$}
\opt{C}{2}
\opt{D}{4}
\end{qoptions}
\qn{3} A uniform slab of mass 7 kg is shown. Find its moment of inertia about the axis normal to the plane of the slab passing through O. Given OA = OC, BC = 10 cm, AB = 2 cm.\par
\qfig{p1_phy_s1_q03_1}
\begin{qoptions}
\opt{A}{84 kg cm$^2$}
\opt{B}{168 kg cm$^2$}
\opt{C}{188 kg cm$^2$}
\opt{D}{198 kg cm$^2$}
\end{qoptions}
\qn{4} An observer at O is at 100 cm from a convex lens of focal length 50 cm. A plane mirror is placed 150 cm behind the lens and starts moving towards the right at 10 cm/s. Find the velocity of her own image as seen by the observer when the mirror just starts moving.\par
\qfig{p1_phy_s1_q04_1}
\begin{qoptions}
\opt{A}{$\dfrac{200}{9}$ cm/s}
\opt{B}{$\dfrac{200}{3}$ cm/s}
\opt{C}{$\dfrac{20}{9}$ cm/s}
\opt{D}{$\dfrac{10}{9}$ cm/s}
\end{qoptions}
\qn{5} A charge distribution is placed on a semicircular rod of radius $R$ with linear charge density $\lambda=\lambda_0\sin\theta$. The potentials at O, P and Q are $V_0$, $V_P$ and $V_Q$ respectively.\par
\qfig{p1_phy_s1_q05_1}
\begin{qoptions}
\opt{A}{$(V_P=V_Q)<V_0$}
\opt{B}{$(V_P=V_Q)>V_0$}
\opt{C}{$V_P=V_Q=V_0$}
\opt{D}{$V_P>V_Q>V_0$}
\end{qoptions}
\qn{6} Radioactive nuclei are produced at a rate $t^2$ and have decay constant $\lambda$. If $\dfrac{dN}{dt}$ is minimum at $t=t_0$, then the number of nuclei $N$ at $t=t_0$ is
\begin{qoptions}
\opt{A}{$\dfrac{2t_0-2\lambda t_0^2}{\lambda^2}$}
\opt{B}{$\dfrac{2t_0-\lambda t_0^2}{\lambda^2}$}
\opt{C}{$\dfrac{t_0-\lambda t_0^2}{\lambda^2}$}
\opt{D}{$\dfrac{2t_0-\lambda t_0^2}{2\lambda^2}$}
\end{qoptions}
\sect{SECTION-I (ii)}
\qn{7} A ring shaped tube of volume $2V_0$ contains an ideal gas with $\gamma=1.5$. It has an adiabatic fixed wall S$_1$ and a movable adiabatic partition S$_2$, and a heater H. S$_2$ moves a quarter of the circle.\par
\qfig{p1_phy_s2_q07_1}
\begin{qoptions}
\opt{A}{Final pressure in chamber 2 is $2P_0$}
\opt{B}{Final pressure in chamber 2 is $2\sqrt2\,P_0$}
\opt{C}{Work done on the gas in chamber 2 by the gas in chamber 1 is $(2\sqrt2-1)P_0V_0$}
\opt{D}{Heat supplied is $4P_0V_0(2\sqrt2-1)$}
\end{qoptions}
\qn{8} A uniformly charged non conducting rod is suspended vertically at its end. The rod can swing freely in the vertical plane without any friction. The linear charge density on the rod is $\lambda$ C/m and it has a uniformly distributed mass of $M$. Length of the rod is $L$. A uniform horizontal electric field ($E$) is switched on in the region. Select the correct option/options :-\par
\qfig{p1_phy_s2_q08_1}
\begin{qoptions}
\opt{A}{If $E=\dfrac{mg}{\lambda L}$, then rod just manages to make it horizontal}
\opt{B}{If $E=\dfrac{mg}{2\lambda L}$, then rod just manages to make it horizontal}
\opt{C}{If $E=\dfrac{Mg}{\lambda L}$, the rod has maximum kinetic energy at an angle $45^\circ$ with vertical}
\opt{D}{If $E=\dfrac{Mg}{2\lambda L}$, the rod has maximum kinetic energy at an angle $45^\circ$ with vertical}
\end{qoptions}
\qn{9} A small ball of mass $m$ is released from rest from a height $h$ above the liquid surface. The density of ball is $\dfrac{\rho}{3}$ : (assume no splash)\par
\qfig{p1_phy_s2_q09_1}
\begin{qoptions}
\opt{A}{If $h\geqslant H$, then the ball will be able to enter the liquid of higher density}
\opt{B}{If $h\geqslant\dfrac{7H}{2}$, then the ball will be able to strike the bottom of the vessel}
\opt{C}{If $h=2H$, then the ball will cross the interface of two liquids with kinetic energy $=mgH$}
\opt{D}{If $h=0$ then ball will float on the surface.}
\end{qoptions}
\qfullbegin
\qn{10} Consider three fixed surfaces shown in the figure. Three blocks each of mass $m$ are released from rest from top of the three surfaces. All blocks reach ground with same speed. Length of path travelled by the blocks is same for second and third surface. If coefficient of friction of three surfaces are $\mu_1$, $\mu_2$ and $\mu_3$ respectively then\par
\qfig{p1_phy_s2_q10_1}
\qfullsplit
\begin{qoptions}
\opt{A}{$\mu_1=\mu_2$}
\opt{B}{$\mu_1>\mu_2$}
\opt{C}{$\mu_2=\mu_3$}
\opt{D}{$\mu_2<\mu_3$}
\end{qoptions}
\qfullend
\qn{11} In the circuit shown in figure $R$ and $V_0$ are constant, while $r$ can change, taking up all values. What can be the equivalent resistance of the circuit ?\par
\qfig{p1_phy_s2_q11_1}
\begin{qoptions}
\opt{A}{1.6 $R$}
\opt{B}{1.4 $R$}
\opt{C}{1.8 $R$}
\opt{D}{2.2 $R$}
\end{qoptions}
\qn{12} There are two inclined planes, as shown in the figure. A ball strikes the inclined plane AB vertically with velocity $v=\sqrt{24g}$ at height $h=\left(1+2\sqrt3\right)$ m. The collision is inelastic. After this the ball moves such that it collides with the inclined plane BC perpendicular to it. Select the option(s).\par
\qfig{p1_phy_s2_q12_1}
\begin{qoptions}
\opt{A}{The co-efficient of restitution between the inclined plane AB and ball is $\dfrac{1}{\sqrt3}$}
\opt{B}{The co-efficient of restitution between the inclined plane AB and the balls is $\dfrac{1}{3\sqrt3}$}
\opt{C}{The time interval between two collision is $\sqrt{\dfrac8g}$ s}
\opt{D}{The velocity of collision with the inclined plane BC is $\sqrt{6g}\left(1+\dfrac{1}{\sqrt3}\right)$ m/s}
\end{qoptions}
\sect{SECTION-II}
\qn{1} A non-viscous incompressible liquid of mass $m$ is filled fully inside a thin uniform spherical shell of mass $m$ and radius $R$ performing pure rolling on a rough horizontal surface. There are two point A and B inside the liquid on the vertical diameter separated by $2R$ as shown in the figure. Pressure difference between B and A is $(P_B-P_A)$. At the given instant velocity of centre of mass of given system is $V_0$ and kinetic energy of this system is $K$. If the value of $\dfrac{K}{P_B-P_A}$ is equal to $\dfrac{n\pi V_0^2R^2}{9g}$. Find $n$.\par
\qfig{p1_phy_s3_q01_1}
\qn{2} A block of mass `m' is given velocity towards another block of mass 6 kg kept at rest on a frictionless horizontal floor some distance apart from a wall as shown in the figure. If all collisions are elastic, find minimum value of mass `m' (in kg) so that blocks will collide only once.\par
\qfig{p1_phy_s3_q02_1}
\qn{3} In a hydrogen like atom an electron is orbiting in an orbit having quantum number $n$. Its frequency of revolution is found to be $13.2\times10^{15}$ Hz. Energy required to free the electron from the atom from the orbit is 54.4 eV. In time 7 nano-seconds the electron jumps back to orbit having quantum number $\dfrac n2$. If $\tau$ be the average torque acted on the electron during the above process, then find $\dfrac{\tau}{3}\times10^{27}$ in N--m. (Given $\dfrac h\pi=2.1\times10^{-34}$ J--s).\par
Frequency of revolution of electron in the ground state of H-atom $f_0=6.6\times10^{15}$ Hz and ionisation energy of H-bond $E_0=13.6$ eV)
\qn{4} As shown in the figure a cylinder of mass $M$ and radius $R$ moves with constant speed $v$ through a region of space that contains dust particles of mass $m$ which are at rest. Number of particles in unit volume is $n$. The cylinder moves in a direction perpendicular to its axis. Assume that the particles do not interact with each other and that $m<<M$. All collisions are elastic and surface of the cylinder is smooth. The force per unit length of the cylinder required to maintain the speed $v$ constant is $\dfrac{K\,nmRv^2}{3}$. Find the value of $K$.\par
\qfig{p1_phy_s3_q04_1}
\qn{5} A positively charged particle is given initial velocity $v_0$ in XY plane as shown in diagram. Initially the particle was at origin and the region contains uniform, mutually perpendicular electric and magnetic field. Let $\vec E=-E\hat j$ and $\vec B=B\hat k$. The minimum magnitude of the magnetic field so that the charge particle crosses the origin again during its subsequent motion is $B_{\min}=\dfrac{n\pi E}{v_0}$. Find $n$\par
\qfig{p1_phy_s3_q05_1}
\qn{6} In the arrangement shown, two springs are attached to a light rod as shown. Length of O to A is $\dfrac{2\ell}{3}$. System is at equilibrium. The time period of oscillation of the block, if it is displaced slightly is found to be $T=2\pi\sqrt{\dfrac{\alpha m}{8K}}$. Find the value of $\alpha$.\par
\qfig{p1_phy_s3_q06_1}
\partsec{2}{CHEMISTRY}
\sect{SECTION-I (i)}
\qn{1} A metal exists in two allotropic forms, $\alpha$ and $\beta$, having a face centred cubic and hexagonal closest packed lattice respectively. Identify the correct statement.
\begin{qoptions}
\opt{A}{Coordination number is higher for the $\alpha$-form.}
\opt{B}{Density is higher for the $\beta$-form.}
\opt{C}{Distance between closest octahedral voids is same for both the forms.}
\opt{D}{Unit cell volume is same for both the forms.}
\end{qoptions}
\qn{2} Consider two Van der waal gases A and B. For gas A, $\displaystyle\lim_{P\to0}\left(\dfrac{\partial Z}{\partial P}\right)_{T=300\text{K}}>0$ and for gas B, $\displaystyle\lim_{P\to0}\left(\dfrac{\partial Z}{\partial P}\right)_{T=300\text{K}}<0$. Identify the incorrect statement.
\begin{qoptions}
\opt{A}{Boyle temperature of gas A $<300$ K}
\opt{B}{Boyle temperature of gas B $>300$K.}
\opt{C}{Boyle temperature of both gases is 300 K}
\opt{D}{Critical temperature of B is higher than that of A}
\end{qoptions}
\qn{3} In which of the following cases, the most stable conformer is achiral?\par
\qfig{p1_chem_s1_q03_1}
\begin{qoptions}
\opt{A}{Only IV}
\opt{B}{I and IV}
\opt{C}{I, II and III}
\opt{D}{In all the cases}
\end{qoptions}
\qfullbegin
\qn{4} \qfig{p1_chem_s1_q04_1}
\qfullsplit
product `D' is:
\begin{qoptions}
\opt{A}{\qfigi{p1_chem_s1_q04_A}}
\opt{B}{\qfigi{p1_chem_s1_q04_B}}
\opt{C}{\qfigi{p1_chem_s1_q04_C}}
\opt{D}{\qfigi{p1_chem_s1_q04_D}}
\end{qoptions}
\qfullend
\qn{5} Select the incorrect statement.
\begin{qoptions}
\opt{A}{The less stable tautomeric form of H$_3$PO$_3$ has three active hydrogens.}
\opt{B}{Both H$_3$PO$_3$ and H$_3$PO$_2$ disproportionate on heating to produce PH$_3$ and H$_3$PO$_4$.}
\opt{C}{In P$_4$O$_{10}$ all the O-P-O bond angles are identical whereas there are two different P-O bond lengths.}
\opt{D}{H$_3$PO$_3$ and H$_3$AsO$_3$ have different basicity.}
\end{qoptions}
\qn{6} Which of the following statement is correct regarding extraction of aluminium from bauxite ore:
\begin{qoptions}
\opt{A}{In Hall's process CO$_2$ is formed at the anode.}
\opt{B}{Baeyer's process is employed when the bauxite ore is rich in silica.}
\opt{C}{In Serpeck's process N$_2$ is produced due to hydrolysis of aluminium nitride.}
\opt{D}{Bauxite ore undergoes roasting before the leaching process.}
\end{qoptions}
\sect{SECTION-I (ii)}
\qn{7} 1 mole of He $\left(C_{v,m}=\dfrac{3R}{2}\right)$ at 1 bar is compressed adiabatically against a constant external pressure, $P_{ext}$ from an initial volume of 20 L. Assuming He to show ideal gas behaviour, identify the correct statement(s) for the process. (Take $R=2$ cal mol$^{-1}$K$^{-1}$, ln3 = 1.1, ln5 = 1.6)\par
\qfig{p1_chem_s2_q07_1}
\begin{qoptions}
\opt{A}{The lowest possible final volume of the gas after this process is 8L.}
\opt{B}{If $P_{ext}=3$ atm, then $\Delta S=0.8$ calK$^{-1}$}
\opt{C}{$\Delta S_{sys}$ is +ve and $\Delta S_{Surroundings}=0$ for all values of $P_{ext}$}
\opt{D}{If $T_1$, $V_1$, are the initial temperature and volume of the system, and $T_2$, $V_2$ are the final temperature and volume respectively, then $\ln\left(\dfrac{T_2}{T_1}\right)=\dfrac{R}{C_v}\ln\left(\dfrac{V_1}{V_2}\right)$}
\end{qoptions}
\qn{8} An azeotrop of HNO$_3$ and H$_2$O is formed with 38 mole percent of HNO$_3$. Identify the correct plot(s).\par
$y\rightarrow$ vapour phase mole fraction.\par
$\chi\rightarrow$ liquid phase mole fraction.
\begin{qoptions}
\opt{A}{\qfigi{p1_chem_s2_q08_A}}
\opt{B}{\qfigi{p1_chem_s2_q08_B}}
\opt{C}{\qfigi{p1_chem_s2_q08_C}}
\opt{D}{\qfigi{p1_chem_s2_q08_D}}
\end{qoptions}
\qfullbegin
\qn{9} \qfig{p1_chem_s2_q09_1}
\qfullsplit
Identify the correct statement(s):
\begin{qoptions}
\opt{A}{Structure of `E' is CH$_3$--CH=CH--CHO.}
\opt{B}{Structure of `B' and `C' can be differentiated by Fehling solution.}
\opt{C}{Structure of `E' is \qfigi{p1_chem_s2_q09_3}}
\opt{D}{Structure of `A' is Ph -- CH = CH -- CH$_3$}
\end{qoptions}
\qfullend
\qn{10} Which of the following is/are not correctly matched?
\begin{qoptions}
\opt{A}{\qfigi{p1_chem_s2_q10_A} $\xrightarrow{\text{DCl}}$ Number of stereoisomers of intermediate and products are identical.}
\opt{B}{\qfigi{p1_chem_s2_q10_B} $\xrightarrow[\text{CCl}_4]{\text{Br}_2}$ Neither product nor intermediate exhibits optical activity.}
\opt{C}{\qfigi{p1_chem_s2_q10_C} $\xrightarrow[\text{Pd}]{\text{H}_2}$ The chair form of the product gives non-superimposable structure after ring flipping.}
\opt{D}{\qfigi{p1_chem_s2_q10_D} $\xrightarrow[\text{PhCOOOCOPh}]{\text{HBr}}$ Product mixture upon fractional distillation gives 4 major fractions, out of which 3 are optically active.}
\end{qoptions}
\qn{11} Select the correct statement(s):
\begin{qoptions}
\opt{A}{[Cr(NH$_3$)$_5$SO$_4$]Cl will produce a white precipitate with AgNO$_3$ solution but not with BaSO$_4$ solution.}
\opt{B}{Octahedral complexes formed by coordinating strong field ligands with Fe$^{+3}$ will be diamagnetic in nature.}
\opt{C}{[Ni(en)$_3$]Cl$_2$, will show optical as well as geometrical isomerism.}
\opt{D}{K[PtCl$_3$C$_2$H$_4$] is an organometallic compound}
\end{qoptions}
\qn{12} A (is a compound of Ba) and B (is a sulphate) are white crystalline substances, highly soluble in water and decompose at higher temperatures. If an aqueous solution of A (slightly basic) is added to an aqueous solution of B (slightly acidic) a white precipitate C forms. The filtrate is neutral and gives a brown colour solution with an excess of acidified KI solution. When boiled, the filtrate evaporates without the appearance of any residue. White solid D can be prepared by the heating of A in the absence of air. The reaction of D with water gives a colourless solution. This solution, if kept in an open container, slowly precipitates a white solid E and leaves water. Upon prolonged exposure to moist air at room temperature, solid D is transformed into E as well. E can be converted back into D but decomposition above 1400\textdegree C is required for this purpose. The reaction of B and D in water forms the precipitate C and is accompanied by a characteristic odour due to formation of a gas. A gives white precipitate with AgNO$_3$. Select the correct option(s):
\begin{qoptions}
\opt{A}{Compound A is on heating gives neutral, acidic and basic oxide}
\opt{B}{Compound B is (NH$_4$)$_2$SO$_4$.}
\opt{C}{Compound E is BaCO$_3$.}
\opt{D}{I$_2$ will precipitate as brown coloured.}
\end{qoptions}
\sect{SECTION-II}
\qn{1} Hydrogen peroxide in aqueous solution decomposes by a first order reaction to produce water and oxygen. The rate constant for this decomposition is $7.50\times10^{-4}$ s$^{-1}$. What quantity of heat (in J) is initially liberated per second from 2.0 L of solution that is 0.02M H$_2$O$_2$? Given $\Delta_fH^\circ$ of H$_2$O$_2$($l$) = $-187.0$ kJ/mol and $\Delta_fH^\circ$ H$_2$O ($l$) = $-287.0$ kJ/mol.
\qn{2} Calculate the pH of solution formed by addition of 1 mole each of HCl, CH$_3$COONa and HCOONa in enough water to make 1 L solution. The pK$_a$ of HCOOH and CH$_3$COOH are 4 and 5, respectively.
\qn{3} In the given reactions:-\par
\qfig{p1_chem_s3_q03_1}
Number of correct statement(s) among the following are :-\par
(1) X and Y are enantiomers\par
(2) M and N are enantiomers\par
(3) A and B are diastereomers\par
(4) (X and Y) are identical with (A and B)\par
(5) (P and Q) are identical with (M and N)\par
(6) M and N are geometrical isomers
\qfullbegin
\qn{4} In how many of the below options, the first structure is more stable than the second.\par
\qfig{p1_chem_s3_q04_1}
\qfullend
\qn{5} If the reaction sequence given below is carried out with 4 mol of BrO$_3^-$ and just required amount of Br$^-$ and H$^+$, the amount of the product D formed (in g) is\underline{\hspace{2cm}}.\par
The yields of A, B, C, D are given in the parentheses.\par
\qfig{p1_chem_s3_q05_1}
\qn{6} The bond angle H-C-H in the formaldehyde molecule is approximately 120\textdegree; the dipole moments for C-H and C = O bonds are $\mu_{C-H}=0.4$ D and $\mu_{C=O}=2.3$ D, respectively. The net dipole moment of the formaldehyde is $X\times10^{-Y}$ Cm (in scientific notation). The value of X is\underline{\hspace{2cm}}.
\partsec{3}{MATHEMATICS}
\sect{SECTION-I (i)}
\qn{1} If $\alpha$ and $\beta$ are exactly two distinct real roots of the equation $(x^2+5x-b^2)(2x^2-ax-6)=0$ where $a,b\in R$, (Given that equation has only real roots) then
\begin{qoptions}
\opt{A}{$\alpha^2+\beta^2=30$}
\opt{B}{$a^2+b^2=105$}
\opt{C}{$\alpha^2+\beta^2=19$}
\opt{D}{$a^2+b^2+\alpha^2+\beta^2=134$}
\end{qoptions}
\qn{2} If $f(x)=x^4-3x^2+7x-1$ and $g(x)=\cot^{-1}(2x^2)+\sin^{-1}\left(\dfrac{1}{1+x^2}\right)$ are two functions, then number of integers in the range of $gof(x)$ is
\begin{qoptions}
\opt{A}{1}
\opt{B}{2}
\opt{C}{3}
\opt{D}{4}
\end{qoptions}
\qn{3} $\displaystyle\lim_{x\to0^+}\dfrac{(1+\tan x)^{2\cot x}-e^{2\cos x}}{1-\cos\left(\ln\left(1+\sin\sqrt x\right)\right)}$
\begin{qoptions}
\opt{A}{$2e^2$}
\opt{B}{$2e$}
\opt{C}{$-2e$}
\opt{D}{$-2e^2$}
\end{qoptions}
\qn{4} $\displaystyle\int e^x\left[\dfrac{2\sin x}{(\sin x+\cos x)^2}-\dfrac{2\cos x}{(\sin x-\cos x)^2}\right]dx$ where `C' is the constant of integration.
\begin{qoptions}
\opt{A}{$e^x\left[\dfrac{2\sin x}{\sin^2x-\cos^2x}\right]+C$}
\opt{B}{$e^x\left[\dfrac{2\cos x}{\sin^2x-\cos^2x}\right]+C$}
\opt{C}{$e^x\left(\dfrac{\sin^2x-\cos^2x}{2\sin x}\right)+C$}
\opt{D}{$e^x\left(\dfrac{\sin^2x-\cos^2x}{2\cos x}\right)+C$}
\end{qoptions}
\qn{5} The value of $S=(\sin^3\theta+\cos^3\theta)+(\sin^33\theta+\cos^33\theta)+(\sin^35\theta+\cos^35\theta)+\dots\dots\dots\dots\dots+(\sin^319\theta+\cos^319\theta)$ when $\theta=\dfrac{\pi}{12}$ is $\dfrac{-a}{\sqrt b}$ where $a\ \&\ b$ are realtively prime numbers, then $|b-6a|$ equals
\begin{qoptions}
\opt{A}{1}
\opt{B}{2}
\opt{C}{3}
\opt{D}{4}
\end{qoptions}
\qn{6} Consider $f(x)=x^3+ax^2+bx+c$ where $a,b,c\in\{1,2,3,4,5,6\}$. The probability that $f(x)=0$ has all real roots given that $f(x)$ is non-decreasing function, is
\begin{qoptions}
\opt{A}{$\dfrac{1}{216}$}
\opt{B}{$\dfrac{1}{96}$}
\opt{C}{$\dfrac{1}{108}$}
\opt{D}{$\dfrac{1}{36}$}
\end{qoptions}
\sect{SECTION-I (ii)}
\qn{7} The vertices of triangle ABC are A(2, 0, 2); B($-1$, 1, 1) and C(1, $-2$, 4). The point D and E divide AB and CA in the ratio 1 : 2 respectively. Another point F is taken in space such that perpendicular drawn from F to the plane containing $\Delta$ABC, meets the plane at the point of intersection (point P) of the line segment CD \& BE. If the distance of F from the plane of triangle ABC is $\sqrt2$ units then
\begin{qoptions}
\opt{A}{$\dfrac{\text{area of }\Delta ABC}{\text{area of }\Delta PBC}=\dfrac{7\sqrt3}{6}$}
\opt{B}{Volume of tetrahedral ABCF $=\dfrac73$ cubic units}
\opt{C}{Equation of line AF is $\vec r=2(\hat i+\hat k)+\lambda(\hat i+2\hat j)$}
\opt{D}{Equation of line AF is $\vec r=2(\hat i+\hat k)+\lambda(\hat i-2\hat k)$}
\end{qoptions}
\qn{8} If chords of the ellipse $\dfrac{x^2}{a^2}+y^2=2$, passing through $(a,-1)$ are bisected by the line $x+y=1$, then
\begin{qoptions}
\opt{A}{least positive integral value of $a$ is 2.}
\opt{B}{least positive integral value of $a$ is 1.}
\opt{C}{sum of integral values of $a$ in $(-10,5)$ is $-8$.}
\opt{D}{sum of negative integral values of $a$ in $(-20,-9)$ is $-145$.}
\end{qoptions}
\qn{9} Let $f:R\longrightarrow R$ be a differentiable function satisfying $f(x+y)=f(x)+f(y)+(x+y)\,e^{x+y}-xe^x-ye^y+4xy$ $\forall$ $x,y\in R$ and $f'(0)=1$, then which of the following option(s) is/are \textbf{correct}?
\begin{qoptions}
\opt{A}{$f(x)=e^{-x+3}$; $x>0$ has exactly one solution}
\opt{B}{$f(2)-f'(2)<0$}
\opt{C}{$3f(x)=\operatorname{sgn}(-x^2+3x-6)$ has exactly one solution in $x\in[2,5]$}
\opt{D}{$f(3)-3e^3=x^2-6x+5$ has exactly two solution}
\end{qoptions}
\qn{10} $\displaystyle\int_0^{\pi}\left[\{x\}+[\cot x]\right]dx$ is greater than or equal to\par
[Note: $\{z\}$ \& $[z]$ denotes fractional part function and greatest integer function of $z$.]
\begin{qoptions}
\opt{A}{0}
\opt{B}{1}
\opt{C}{$-\dfrac\pi2$}
\opt{D}{$-\pi$}
\end{qoptions}
\qn{11} Let OABC is a rhombus where O is origin such that vertices A ($z_1$) and C ($z_3$) lie on the circle $|z|=2$ and the vertex B ($z_2$) lie on the circle $|z|=3$. Let Arg $(z_2-z_1)=\dfrac\pi3$, then
\begin{qoptions}
\opt{A}{$z_1=\left(\dfrac{3-\sqrt7 i}{3+\sqrt7 i}\right)(1+\sqrt3 i)$}
\opt{B}{$z_2=\left(\dfrac{6+6\sqrt3 i}{3-\sqrt7 i}\right)$}
\opt{C}{$z_3=(1+\sqrt3 i)$}
\opt{D}{$z_1=\left(\dfrac{3-\sqrt7 i}{3+\sqrt7 i}\right)(1-\sqrt3 i)$}
\end{qoptions}
\qn{12} Let $f(x)=-x^3+px^2+qx+6\operatorname{sgn}(x^2+2x+5)$ where $p,q\in R$ \& sgn denotes signum function. If the largest possible interval in which $f'(x)$ is positive is $(-2,1)$, then $(p+q)$ is less than or equal to
\begin{qoptions}
\opt{A}{4}
\opt{B}{5}
\opt{C}{6}
\opt{D}{9}
\end{qoptions}
\sect{SECTION-II}
\qn{1} Let a plane P containing a line $\vec r=\left(\hat i+\hat j+\hat k\right)+\lambda\left(-\hat i+3\hat j+2\hat k\right)$ and perpendicular to $2x+3y-z=2$. If the area of projection triangle formed on plane P whose vertices are $A_r=2r\hat i+3(-1)^r\hat j+r^2\hat k$ where $r\in N\ \&\ 1\leqslant r\leqslant3$ is $a\sqrt b$ where $a, b$ are relatively prime numbers, then $\left(\dfrac ba\right)$ is equal to
\qn{2} In a book with page numbers 1 to $n$, four consecutive pages are torn off and the sum of the numbers on the remaining pages is 1596, then find the sum of possible values of $n$. (4 consecutive pages means page numbers $m$, $m+1$, $m+2$ \& $m+3$, $m\in N$ )
\qn{3} From a fixed point A, three normals are drawn to the parabola $y^2=4ax$ at the points P, Q and R. Two circles $C_1$ and $C_2$ be drawn on AP and AQ as diameter. If slope of the common chord of the circles $C_1$ and $C_2$ be $m_1$ and the slope of the tangent to the parabola at the point R be $m_2$, then $m_1\times m_2$ is equal to
\qn{4} Let S be the set containing all matrices $A=[a_{ij}]_{3\times3}$ having entries from the set $\{0,1,2,4\}$ such that sum of all principal diagonal entries of $AA^T=17$ then number of matrices in S is
\qn{5} If $\displaystyle\int_0^1\left((6x^4-f(x))\right)f(x)\,dx=1$, then the area of region enclosed by $y=f(x)$, $x=0$ and $x=2$.
\qn{6} Let $u=\sin^7x$, $v=\cos x$. If $\dfrac{d^2}{dx^2}\left(\dfrac{d^3u}{dv^3}\right)\Bigg|_{x=\frac\pi8}=2^{\lambda_1}\cdot3^{\lambda_2}\cdot5^{\lambda_3}\cdot7^{\lambda_4}$, then the value of $(\lambda_1+\lambda_2+\lambda_3+\lambda_4)$, is
\end{multicols}
\end{document}
